\chapter{Validazione software e simulazione}

Questa sezione descrive come validare il comportamento del firmware in simulazione e su ambiente di test, senza entrare in aspetti elettrici.

\section{Strategia di validazione}
\begin{itemize}
  \item \textbf{Unit test}: coprono le funzioni di base e la logica dei moduli.
  \item \textbf{Simulazione Wokwi}: verifica end-to-end del ciclo \texttt{loop()} con sensori e attuatori virtuali.
  \item \textbf{Test di regressione}: esecuzione ripetuta della suite dopo modifiche al codice.
\end{itemize}

\section{Procedure di validazione in simulazione}
\begin{enumerate}
  \item Avviare la simulazione e verificare l'inizializzazione del display.
  \item Premere i pulsanti IN/OUT e verificare conteggio persone e semaforo.
  \item Forzare temperature alte/basse e osservare LED di heating/cooling.
  \item Modificare umidita e osservare attivazione del controllo.
  \item Oscurare o illuminare la fotoresistenza e verificare la plafoniera.
  \item Controllare la rotazione delle pagine LCD e i messaggi di errore.
\end{enumerate}

\section{Criteri di accettazione software}
\begin{itemize}
  \item Nessun crash o blocco nel ciclo \texttt{loop()}.
  \item Tutti i moduli rispondono coerentemente ai valori di input simulati.
  \item Le condizioni di errore disattivano le uscite e mostrano diagnostica.
  \item La suite di test e eseguita con esito positivo.
\end{itemize}
